% Generated by book/tools/notebook_to_tex.py. Do not edit by hand.

\chapter[다중 클래스 분류 (Multi-class Classification \& CE)]{다중 클래스 분류 (Multi-class Classification \& CE)}

\label{ch:05}

\markboth{다중 클래스 분류 (Multi-class Classification \& CE)}{다중 클래스 분류 (Multi-class Classification \& CE)}

\chaptermeta{05_sklearn_multiclass/05_sklearn_multiclass.ipynb}{https://colab.research.google.com/github/yoon-gu/neuqes-101/blob/master/05_sklearn_multiclass/05_sklearn_multiclass.ipynb}{K=5 출력 헤드와 softmax 일반화}



\index{1Multi class classification@Multi-class classification}

\index{1CrossEntropyLoss@CrossEntropyLoss}

\index{1confusion matrix@confusion\_matrix}

\index{1classification report@classification\_report}

\index{0다중 클래스 분류@다중 클래스 분류}

\index{0균등 추측@균등 추측}

\index{0클래스 불균형@클래스 불균형}

\index{0혼동 행렬@혼동 행렬}

\index{0매크로 F1@매크로 F1}

\index{1multinomial logistic regression@multinomial logistic regression}

\index{1OvR@OvR}

\index{1One vs Rest@One-vs-Rest}

\index{1argmax@argmax}

\index{1weighted F1@weighted F1}

\index{1macro average@macro average}

\index{1weighted average@weighted average}

\index{1baseline@baseline}

\index{1log K@log K}

\index{1precision@precision}

\index{1recall@recall}

\index{1F1@F1}

\index{1class weight@class\_weight}

\index{1multi class@multi\_class}

\index{0다항 로지스틱 회귀@다항 로지스틱 회귀}

\index{0일대나머지@일대나머지}

\index{0상호배타 클래스@상호배타 클래스}

\index{0클래스 경쟁@클래스 경쟁}

\index{0가중 F1@가중 F1}

\index{0매크로 평균@매크로 평균}

\index{0가중 평균@가중 평균}

\index{0기준선@기준선}

\index{0분류 리포트@분류 리포트}



\section{변화추적표}

\begin{booktable}{5장 변화추적표}{tab:ch05-01}
\begin{adjustbox}{max width=\textwidth}
\begin{tabular}{@{}lllllll@{}}
\toprule
Ch & 모델 & 토크나이저 & 데이터 & Output Head & Activation & Loss \\
\midrule

1 & (TF-IDF) & \inlinecode{TfidfVectorizer()} & Yelp 5,000 & --- & --- & --- \\
2 & \inlinecode{LinearRegression()} & \inlinecode{TfidfVectorizer()} & Yelp (별점 1-5) & (1차원) & 없음 & \inlinecode{MSELoss} \\
3 & \inlinecode{LogisticRegression()} & \inlinecode{TfidfVectorizer()} & Yelp 이진화 & (1차원) & sigmoid & \inlinecode{BCEWithLogitsLoss} \\
4 & \inlinecode{LogisticRegression()} (multinomial 자동) & \inlinecode{TfidfVectorizer()} & Yelp 이진화 (\ref{ch:03}장과 동일) & (2차원) & softmax & \inlinecode{CrossEntropyLoss} \\
\textbf{5 \(\leftarrow\) 여기} & \inlinecode{LogisticRegression()} (multinomial 자동) & \inlinecode{TfidfVectorizer()} & Yelp 5클래스 (별점 0-4) & \textbf{(5차원)} & softmax & \inlinecode{CrossEntropyLoss} \\
\bottomrule
\end{tabular}
\end{adjustbox}
\end{booktable}

전체 \ref{ch:20}장 표는 \href{https://github.com/yoon-gu/neuqes-101\#챕터별-변화추적표}{루트 README.md}를 참고하세요.



\section{변경점: \ref{ch:04}장 대비}

\begin{booktable}{5장 변경점 요약}{tab:ch05-02}
\begin{adjustbox}{max width=\textwidth}
\begin{tabular}{@{}lll@{}}
\toprule
축 & \ref{ch:04}장 & \ref{ch:05}장 \\
\midrule

데이터 & Yelp 이진화 (K=2) & \textbf{Yelp 5클래스 (K=5)} \\
Output Head & (2차원) & \textbf{(5차원)} \\
활성화·Loss & softmax + CE & softmax + CE (그대로) \\
모델·토크나이저 & LogReg(multinomial) + TF-IDF & 그대로 \\
\bottomrule
\end{tabular}
\end{adjustbox}
\end{booktable}

\textbf{한 가지 변화} --- K=2가 K=5로 늘어난 것. softmax/CE는 K가 무엇이든 같은 식이라 코드 변화도 거의 없습니다. 이게 softmax/CE의 진짜 가치 --- sigmoid+BCE는 K=2 전용이지만 softmax+CE는 자연스럽게 다중 클래스로 확장됩니다.



\section{손실 노트 --- 같은 교차 엔트로피, K=5 수치 예시}

Loss는 \ref{ch:04}장와 동일한 \inlinecode{CrossEntropyLoss}. K가 늘어나도 식은 그대로:

\begin{equation}
\label{eq:ch05-01}
L = -\frac{1}{N}\sum_{i=1}^{N} \log \hat p_{i,\, y_i}
\end{equation}

식~\eqref{eq:ch05-01}은 이 절에서 사용하는 핵심 관계를 정리한 것입니다.

\textbf{숫자로 감 잡기} (K=5, 정답 클래스 = 2 인 한 샘플):

\begin{booktable}{5장 손실 수치 예시}{tab:ch05-03}
\begin{adjustbox}{max width=\textwidth}
\begin{tabular}{@{}lll@{}}
\toprule
예측 분포 \([\hat p_0, \hat p_1, \hat p_2, \hat p_3, \hat p_4]\) & 정답 확률 \(\hat p_2\) & 손실 \(-\log \hat p_2\) \\
\midrule

\textbf{정답에 집중}: \texttt{{[}0.05,\ 0.05,\ 0.80,\ 0.05,\ 0.05{]}} & 0.80 & 0.223 \\
\textbf{균등(uniform)}: \texttt{{[}0.20,\ 0.20,\ 0.20,\ 0.20,\ 0.20{]}} & 0.20 & 1.609 \\
\textbf{틀린 클래스에 집중}: \texttt{{[}0.05,\ 0.05,\ 0.05,\ 0.05,\ 0.80{]}} & 0.05 & \textbf{2.996} \\
\bottomrule
\end{tabular}
\end{adjustbox}
\end{booktable}

\textbf{baseline = \(\log K = \log 5 \approx 1.609\)}: 모델이 아무 정보 없이 균등 추측만 할 때의 손실. 학습된 모델은 이보다 작아야 하고, baseline을 초과하면 “정답이 \emph{아닌} 곳에 자신 있다”는 신호 --- gradient가 모델을 강하게 끌어당깁니다.



\section{토크나이저 노트}

\textbf{\ref{ch:01}--\ref{ch:04}장와 동일한 \inlinecode{TfidfVectorizer}}. 입력 표현, 모델, loss 모두 그대로 --- 변하는 건 데이터의 클래스 수와 출력 차원뿐.

\begin{previewBox}{미리보기}
\textbf{다음 장(\ref{ch:06}장)}: 같은 TF-IDF. 변화는 활성화 함수가 softmax(합=1)에서 \textbf{K개 독립 sigmoid}(라벨 간 독립)로 일반화되어 multi-label로 확장.
\end{previewBox}



\Needspace{5\baselineskip}
\begin{lstlisting}
!pip install -q datasets scikit-learn pandas matplotlib
\end{lstlisting}

\begin{codeRead}
\textbf{1행}의 \inlinecode{!pip install -q datasets...}에서는 Colab 실행에 필요한 패키지를 설치합니다.
\end{codeRead}



\Needspace{11\baselineskip}
\begin{lstlisting}
import warnings
warnings.filterwarnings("ignore", category=FutureWarning)
warnings.filterwarnings(
    "ignore",
    category=DeprecationWarning,
)

import numpy as np
import pandas as pd
import matplotlib.pyplot as plt
from datasets import load_dataset
from sklearn.feature_extraction.text import TfidfVectorizer
from sklearn.linear_model import LogisticRegression
from sklearn.multiclass import OneVsRestClassifier
from sklearn.model_selection import train_test_split
from sklearn.metrics import accuracy_score, classification_report, confusion_matrix
\end{lstlisting}

\noindent\emph{앞 블록은 필요한 라이브러리와 기본 설정을 준비하는 단계입니다. 이어지는 블록에서는 데이터를 불러오고 학습에 맞는 형태로 정리하는 단계로 넘어갑니다.}

\Needspace{11\baselineskip}
\begin{lstlisting}[firstnumber=17]
plt.rcParams["axes.unicode_minus"] = False

dataset = load_dataset("Yelp/yelp_review_full")
SAMPLE_SIZE = 5000
ds = dataset["train"].shuffle(seed=42).select(range(SAMPLE_SIZE))
df = ds.to_pandas()
print(f"Total samples: {len(df)}")
print("Class distribution (label 0-4 = star 1-5):")
print(df["label"].value_counts().sort_index())
\end{lstlisting}

\begin{codeRead}
\inlinecode{numpy, pandas, matplotlib, datasets, scikit-learn} 같은 주요 패키지를 불러옵니다.
\end{codeRead}

\noindent\textbf{출력.}
\begin{lstlisting}[style=bookoutput]
If the error persists, please let us know by opening an issue on GitHub (ht...
Total samples: 5000
Class distribution (label 0-4 = star 1-5):
label
0    1017
1    1027
2     960
3    1021
4     975
Name: count, dtype: int64
\end{lstlisting}
\par\vspace{0.9em}



\Needspace{11\baselineskip}
\begin{lstlisting}
y_5class = df["label"]

X_text_train, X_text_test, y_train, y_test = train_test_split(
    df["text"], y_5class, test_size=0.2, random_state=42, stratify=y_5class,
)

tfidf = TfidfVectorizer(max_features=10000)
X_train = tfidf.fit_transform(X_text_train)
X_test = tfidf.transform(X_text_test)

print(
    f"X_train: {X_train.shape}, y_train distribution: "
    f"{pd.Series(y_train).value_counts().sort_index().tolist()}"
)
\end{lstlisting}

\begin{codeRead}
\textbf{1행}의 \inlinecode{y\_5class = ...}에서는 이후 분석에서 사용할 값을 준비합니다. \textbf{3--5행}의 \inlinecode{X\_text\_train, X\_text\_test...}에서는 학습용 데이터와 평가용 데이터를 분리합니다. \textbf{7행}의 \inlinecode{tfidf = ...}에서는 TF-IDF 기반 벡터라이저를 만듭니다. \textbf{8행}의 \inlinecode{X\_train = ...}에서는 텍스트나 라벨을 모델이 처리할 수 있는 수치 표현으로 변환합니다. \textbf{9행}의 \inlinecode{X\_test = ...}에서는 텍스트나 라벨을 모델이 처리할 수 있는 수치 표현으로 변환합니다.
\end{codeRead}

\noindent\textbf{출력.}
\begin{lstlisting}[style=bookoutput]
X_train: (4000, 10000), y_train distribution: [813, 822, 768, 817, 780]
\end{lstlisting}
\par\vspace{0.9em}



\section{실습: 5클래스 분류}

코드는 \ref{ch:04}장과 동일한 \inlinecode{LogisticRegression()} 한 줄. sklearn 은 \emph{학습 데이터의 클래스 개수} 를 보고, K=5 multi-class 라 자동으로 multinomial(softmax+CE)을 적용합니다.



\Needspace{11\baselineskip}
\begin{lstlisting}
model_5 = LogisticRegression(max_iter=1000)
model_5.fit(X_train, y_train)

y_pred = model_5.predict(X_test)
acc = accuracy_score(y_test, y_pred)
baseline = 1 / 5

print(f"Test accuracy: {acc:.4f}")
print(f"baseline (uniform guess): {baseline:.4f}")
print(f"Improvement over baseline: {acc - baseline:+.4f}")
\end{lstlisting}

\begin{codeRead}
\textbf{1행}의 \inlinecode{model\_5 = ...}에서는 이번 실습에서 관찰할 모델 객체를 정의합니다. \textbf{2행}의 \inlinecode{model\_5.fit(...)}에서는 학습 데이터로 모델 또는 변환기를 적합합니다. \textbf{4행}의 \inlinecode{y\_pred = ...}에서는 학습된 모델로 최종 예측값을 만듭니다. \textbf{5행}의 \inlinecode{acc = ...}에서는 이후 분석에서 사용할 값을 준비합니다. \textbf{6행}의 \inlinecode{baseline = ...}에서는 이후 분석에서 사용할 값을 준비합니다.
\end{codeRead}

\noindent\textbf{출력.}
\begin{lstlisting}[style=bookoutput]
Test accuracy: 0.5110
baseline (uniform guess): 0.2000
Improvement over baseline: +0.3110
\end{lstlisting}
\par\vspace{0.9em}



\Needspace{11\baselineskip}
\begin{lstlisting}
proba_5 = model_5.predict_proba(X_test)
print(f"predict_proba shape: {proba_5.shape}  (N, K=5)")
print(
    f"Row sums (should be 1): "
    f"{proba_5.sum(axis=1)[:5].round(4)}"
)
print(f"\nFirst 3 sample probability distributions:")
print(
    pd.DataFrame(proba_5[:3], columns=[f"P({i+1}★)" for i in range(5)]).round(3)
)
\end{lstlisting}

\begin{codeRead}
\textbf{1행}의 \inlinecode{proba\_5 = ...}에서는 클래스별 예측 확률을 계산합니다.
\end{codeRead}

\noindent\textbf{출력.}
\begin{lstlisting}[style=bookoutput]
predict_proba shape: (1000, 5)  (N, K=5)
Row sums (should be 1): [1. 1. 1. 1. 1.]

First 3 sample probability distributions:
   P(1★)  P(2★)  P(3★)  P(4★)  P(5★)
0  0.357  0.307  0.163  0.081  0.092
1  0.136  0.339  0.348  0.138  0.040
2  0.614  0.233  0.052  0.042  0.059
\end{lstlisting}
\par\vspace{0.9em}



\section{해부: 평가 지표 한꺼번에 보기}

5클래스에서는 클래스마다 precision/recall/F1이 따로 정의됩니다. \inlinecode{classification\_report}가 한 번에 정리.



\Needspace{5\baselineskip}
\begin{lstlisting}
print(
    classification_report(y_test, y_pred, target_names=[f"{i+1}★" for i in range(5)])
)
\end{lstlisting}

\noindent\textbf{출력.}
\begin{lstlisting}[style=bookoutput]
precision    recall  f1-score   support

          1★       0.62      0.77      0.69       204
          2★       0.40      0.35      0.37       205
          3★       0.41      0.35      0.38       192
          4★       0.45      0.43      0.44       204
          5★       0.62      0.64      0.63       195

    accuracy                           0.51      1000
   macro avg       0.50      0.51      0.50      1000
weighted avg       0.50      0.51      0.50      1000
\end{lstlisting}
\par\vspace{0.9em}



\Needspace{9\baselineskip}
\begin{lstlisting}
cm = confusion_matrix(y_test, y_pred)
cm_df = pd.DataFrame(
    cm,
    index=[f"true {i+1}★" for i in range(5)],
    columns=[f"pred {i+1}★" for i in range(5)],
)
print(cm_df)
\end{lstlisting}

\begin{codeRead}
\textbf{1행}의 \inlinecode{cm = ...}에서는 이후 분석에서 사용할 값을 준비합니다. \textbf{2--6행}의 \inlinecode{cm\_df = ...}에서는 결과를 표로 보기 좋게 정리합니다.
\end{codeRead}

\noindent\textbf{출력.}
\begin{lstlisting}[style=bookoutput]
pred 1★  pred 2★  pred 3★  pred 4★  pred 5★
true 1★      158       36        2        4        4
true 2★       54       72       52       16       11
true 3★       23       50       68       37       14
true 4★       11       20       37       88       48
true 5★       10        3        6       51      125
\end{lstlisting}
\par\vspace{0.9em}



\textbf{관찰}: confusion matrix의 오답이 \textbf{대각선 근처에 몰립니다}.

\begin{itemize}
\tightlist
\item
  4점을 5점으로 혼동하는 경우는 많아도 4점을 1점으로 혼동하는 경우는 거의 없습니다.
\item
  분류 모델은 클래스 간 거리(별점 1과 별점 2가 가깝다는 정보)를 \emph{명시적으로} 모릅니다 --- 5개 독립 클래스로 처리할 뿐입니다.
\item
  그런데도 오답이 대각선 근처에 모이는 이유: \textbf{데이터 자체가 ordinal}. 1점 리뷰의 단어 분포와 2점 리뷰의 단어 분포가 비슷해서, 모델이 자연스럽게 “비슷한 클래스끼리 혼동” 패턴을 학습.
\end{itemize}

\textbf{회귀(\ref{ch:02}장)와 비교} --- 같은 데이터, 다른 관점:

\begin{booktable}{5장 해부: 평가 지표 한꺼번에 보기 표}{tab:ch05-04}
\begin{adjustbox}{max width=\textwidth}
\begin{tabular}{@{}llll@{}}
\toprule
관점 & 가정 & 손실 & 4점을 5점으로 vs 1점으로 \\
\midrule

회귀 (\ref{ch:02}장) & 별점 사이 \textbf{거리} 가 의미 있음 & MSE & 1 vs 16 (다른 페널티) \\
분류 (\ref{ch:05}장) & 5개 \textbf{독립 클래스} & CE & 둘 다 동일한 오분류 (같은 페널티) \\
\bottomrule
\end{tabular}
\end{adjustbox}
\end{booktable}

회귀는 ordinal 정보를 loss에 반영하고, 분류는 클래스별 패턴(예: 1점 리뷰 욕설 vs 5점 리뷰 칭찬)에 더 집중합니다. 둘 중 어느 게 나은지는 케바케.



\section{변형: multinomial vs OvR}

\textbf{multinomial} (\inlinecode{LogisticRegression()} 의 모던 sklearn 기본 동작) 은 한 모델이 K개 logit을 \textbf{동시에} 학습합니다. softmax 한 번이라 합 = 1이 강제 --- “K개 클래스 중 정확히 하나”라는 \emph{상호배타} 가정.

또 다른 방식 \textbf{OvR (One-vs-Rest)} 은 K개의 \emph{독립} binary 분류기. 각 분류기는 “이 클래스 vs 나머지 모든 클래스”만 학습합니다.

\subsection{두 방식의 구조 비교}

\textbf{multinomial (softmax)}:

\Needspace{8\baselineskip}
\begin{lstlisting}
[입력 x] ──$\to$ Linear(V $\to$ K) ──$\to$ logits [z_1, ..., z_K]
                            ──$\to$ softmax 한 번
                            ──$\to$ [p_1, ..., p_K]   (
                                합 = 1,
                                클래스끼리 경쟁,
                            )
\end{lstlisting}

\textbf{OvR} (이 장의 K=5 예시):

\Needspace{10\baselineskip}
\begin{lstlisting}
             ┌──$\to$ 분류기 1:  "1★ vs 나머지"  ──$\to$ sigmoid ──$\to$ P_1
             ├──$\to$ 분류기 2:  "2★ vs 나머지"  ──$\to$ sigmoid ──$\to$ P_2
[입력 x] ──$\to$ ├──$\to$ 분류기 3:  "3★ vs 나머지"  ──$\to$ sigmoid ──$\to$ P_3
             ├──$\to$ 분류기 4:  "4★ vs 나머지"  ──$\to$ sigmoid ──$\to$ P_4
             └──$\to$ 분류기 5:  "5★ vs 나머지"  ──$\to$ sigmoid ──$\to$ P_5

   각 P_k 는 다른 P_j 와 무관한 독립 sigmoid 출력 (raw 합 $\ne$ 1)
   예측: argmax(P_k) — sklearn이 표시할 땐 행을 정규화해 합 1로 보여줌
\end{lstlisting}

핵심 차이는 \textbf{클래스가 서로 경쟁하느냐} 입니다. multinomial은 한 logit이 커지면 다른 logit의 softmax 확률이 자동으로 줄어듭니다 (분모 공유). OvR의 각 sigmoid는 다른 클래스 학습과 독립적이라 P\_k 가 모두 동시에 0.8이 될 수도, 모두 0.1이 될 수도 있습니다.



\Needspace{11\baselineskip}
\begin{lstlisting}
model_ovr = OneVsRestClassifier(LogisticRegression(max_iter=1000))
model_ovr.fit(X_train, y_train)

print(
    f"OvR estimators count: "
    f"{len(model_ovr.estimators_)}"
)
print(
    f"Each estimator is a separate LogisticRegression for"
    f"'class k vs rest'"
)
print(
    f"  estimator 0 coef_ shape: "
    f"{model_ovr.estimators_[0].coef_.shape}  (1, V)"
)

\end{lstlisting}

\noindent\emph{앞 블록은 앞에서 만든 중간 값을 다음 계산으로 넘기는 단계입니다. 이어지는 블록에서는 모델 출력을 확률이나 예측값으로 바꾸는 단계로 넘어갑니다.}

\Needspace{11\baselineskip}
\begin{lstlisting}[firstnumber=17]
ovr_coef = np.vstack([est.coef_[0] for est in model_ovr.estimators_])
# (5, V)
ovr_intercept = np.array([est.intercept_[0] for est in model_ovr.estimators_]) # (5,)

ovr_logits_all = np.asarray(X_test @ ovr_coef.T) + ovr_intercept
ovr_sigmoid_all = 1.0 / (1.0 + np.exp(-ovr_logits_all))

sample_idx = 0
sample_text = X_text_test.iloc[sample_idx]
true_label = y_test.iloc[sample_idx]

p_multi = proba_5[sample_idx]
p_ovr_raw = ovr_sigmoid_all[sample_idx]
p_ovr_norm = model_ovr.predict_proba(X_test)[sample_idx]

\end{lstlisting}

\noindent\emph{앞 블록은 모델 출력을 확률이나 예측값으로 바꾸는 단계입니다. 이어지는 블록에서는 앞에서 만든 중간 값을 다음 계산으로 넘기는 단계로 넘어갑니다.}

\Needspace{11\baselineskip}
\begin{lstlisting}[firstnumber=32]
print("\nReview preview (200 chars):")
print(f"{sample_text[:200]}...")
print(f"True star:    {true_label + 1}★\n")

print(
    f"{'class':>8}  {'multinomial':>14}  {'OvR raw':>10}  "
    f"{'OvR normalized':>16}"
)
print("-" * 56)
\end{lstlisting}

\noindent\emph{앞뒤 블록은 모두 앞에서 만든 중간 값을 다음 계산으로 넘기는 단계입니다. 길이를 나누어 같은 흐름을 단계별로 확인합니다.}

\Needspace{11\baselineskip}
\begin{lstlisting}[firstnumber=41]
for k in range(5):
    print(
        f"  {k+1}★    {p_multi[k]:>14.4f}  {p_ovr_raw[k]:>10.4f}  "
        f"{p_ovr_norm[k]:>16.4f}"
    )
print("-" * 56)
print(
    f"  sum    {p_multi.sum():>14.4f}  {p_ovr_raw.sum():>10.4f}  "
    f"{p_ovr_norm.sum():>16.4f}"
)
\end{lstlisting}

\noindent\textbf{출력.}
\begin{lstlisting}[style=bookoutput]
OvR estimators count: 5
Each estimator is a separate LogisticRegression for 'class k vs rest'
  estimator 0 coef_ shape: (1, 10000)  (1, V)

Review preview (200 chars):
A friend suggested this.. and we met up here with 2 other friends on Christ...
True star:    2★

   class     multinomial     OvR raw    OvR normalized
--------------------------------------------------------
  1★            0.3570      0.3374            0.3625
  2★            0.3071      0.2680            0.2880
  3★            0.1634      0.1543            0.1658
  4★            0.0806      0.0827            0.0889
  5★            0.0919      0.0882            0.0948
--------------------------------------------------------
  sum            1.0000      0.9307            1.0000
\end{lstlisting}
\par\vspace{0.9em}



\textbf{관찰}

\begin{itemize}
\tightlist
\item
  \textbf{multinomial 열}: 깨끗한 분포, 합 = 1. “이 문서는 K개 별점 중 어느 \emph{한} 별점일 확률”을 나타냄.
\item
  \textbf{OvR raw 열}: 5개 sigmoid가 서로 독립적으로 작동한 결과. 합이 1이 아닙니다 (보통 1보다 크거나 작음).
\item
  \textbf{OvR 정규화 후 열}: sklearn이 raw 값을 행 합으로 나눠 합=1을 만들어준 것. 표시용일 뿐 모델의 본래 출력은 아닙니다.
\end{itemize}

\textbf{왜 이 차이가 중요한가} (\ref{ch:06}장 후속 논점)

\begin{itemize}
\tightlist
\item
  multinomial의 합=1 제약은 “이 문서의 별점은 정확히 \emph{하나} 다”라는 데이터 가정에 잘 맞습니다.
\item
  그러나 한 문서가 여러 라벨을 가질 수 있는 \emph{multi-label} 문제에서는 이 가정이 깨집니다 --- 영화는 “로맨스 + 코미디”일 수 있고, 식당 리뷰는 “음식 + 서비스 + 가격”을 동시에 다룰 수 있습니다.
\item
  multi-label은 \textbf{OvR의 사고방식을 그대로} 가져갑니다: K개 독립 sigmoid를 별도로 학습하고, \textbf{정규화하지 않습니다}. 각 라벨이 독립적으로 0/1을 결정하는 것이 그 모델의 본래 모습.
\item
  그래서 OvR을 multi-class에서 미리 만나두는 게 다음 장의 다리가 됩니다.
\end{itemize}



\Needspace{11\baselineskip}
\begin{lstlisting}
acc_ovr = accuracy_score(
    y_test,
    model_ovr.predict(X_test),
)
print(f"multinomial accuracy: {acc:.4f}")
print(f"OvR accuracy:         {acc_ovr:.4f}")
print(f"Diff: {abs(acc - acc_ovr):.4f}")

raw_sums = ovr_sigmoid_all.sum(axis=1)
print(
    f"\nOvR raw row sum distribution (pre-normalization):"
)
print(f"  min:  {raw_sums.min():.3f}")
print(f"  max:  {raw_sums.max():.3f}")
print(f"  mean: {raw_sums.mean():.3f}")
print(
\end{lstlisting}

\noindent\emph{앞 블록은 모델 출력을 확률이나 예측값으로 바꾸는 단계입니다. 이어지는 블록에서는 앞에서 만든 중간 값을 다음 계산으로 넘기는 단계로 넘어갑니다.}

\Needspace{5\baselineskip}
\begin{lstlisting}[firstnumber=17]
    f"  -> 5 independent sigmoids; rows do not sum to "
    f"exactly 1"
)
\end{lstlisting}

\noindent\textbf{출력.}
\begin{lstlisting}[style=bookoutput]
multinomial accuracy: 0.5110
OvR accuracy:         0.5070
Diff: 0.0040

OvR raw row sum distribution (pre-normalization):
  min:  0.778
  max:  1.387
  mean: 0.982
  → 5 independent sigmoids; rows do not sum to exactly 1
\end{lstlisting}
\par\vspace{0.9em}



\section{이 장에 등장한 라이브러리}

\begin{booktable}{5장 새로 등장한 라이브러리}{tab:ch05-05}
\begin{adjustbox}{max width=\textwidth}
\begin{tabular}{@{}lll@{}}
\toprule
이름 & 한 줄 설명 & 다음 장에서 \\
\midrule

\inlinecode{sklearn.multiclass.OneVsRestClassifier} & K개 독립 binary 분류기를 묶어 estimators\_ 로 노출 & \ref{ch:06}장 multi-label의 핵심 도구로 그대로 재등장 \\
\inlinecode{sklearn.metrics.confusion\_matrix} & 다중 클래스도 K\(\times\)K로 일반화 & 분류 장마다 사용 \\
\inlinecode{sklearn.metrics.classification\_report} & per-class precision/recall/F1 한 번에 & --- \\
\bottomrule
\end{tabular}
\end{adjustbox}
\end{booktable}



\section{체크포인트 질문}

\begin{enumerate}
\def\labelenumi{\arabic{enumi}.}
\tightlist
\item
  K=5 균등 분포일 때 CE 손실값은 얼마인가요? 학습된 모델은 이보다 작아야 하는 이유는?
\item
  5\(\times\)5 confusion matrix에서 오답이 대각선 근처에 몰리는 현상은 모델의 어떤 가정과 데이터의 어떤 성질 사이에서 나오나요?
\item
  multinomial과 OvR의 핵심 차이를 한 문장으로 요약해보세요.
\item
  같은 별점 데이터를 회귀(\ref{ch:02}장)로 풀 때와 5클래스 분류(\ref{ch:05}장)로 풀 때 어떤 종류의 실수에 더 큰 페널티를 주나요?
\end{enumerate}



\section{FAQ}

\begin{faqBox}{Q1. (이론) multinomial과 OvR은 어떻게 다른가요?}

핵심 차이는 \textbf{클래스 간 의존성}.

\begin{itemize}
\tightlist
\item
  \textbf{multinomial (softmax)}: 한 모델이 K개 logit을 동시에 학습. softmax \(\to\) 확률 합 = 1, 클래스 \emph{상호배타}.
\item
  \textbf{OvR (One-vs-Rest)}: K개 \emph{독립} binary 분류기 (\inlinecode{OneVsRestClassifier}로 명시적으로 구성). 클래스 0의 logit은 다른 클래스 학습에 영향 없음. 정규화 전 확률 합이 1이 아닐 수 있음.
\end{itemize}

\begin{booktable}{5장 FAQ 표}{tab:ch05-06}
\begin{adjustbox}{max width=\textwidth}
\begin{tabular}{@{}ll@{}}
\toprule
상황 & 적합한 방식 \\
\midrule

한 샘플에 정확히 한 라벨 (별점, 뉴스 카테고리) & multinomial \\
한 샘플에 여러 라벨 가능 (영화 장르: 로맨스+코미디) & OvR (\ref{ch:06}장) \\
클래스 수백 개 + 빠른 학습 필요 & OvR (binary들이 병렬 학습 쉬움) \\
\bottomrule
\end{tabular}
\end{adjustbox}
\end{booktable}

\end{faqBox}

\begin{faqBox}{Q2. (실무) 클래스 수가 100개를 넘어가면 학습이 느려지는데 어떻게 하나요?}

세 가지 흔한 처리법.

\begin{enumerate}
\def\labelenumi{\arabic{enumi}.}
\tightlist
\item
  \textbf{Hierarchical classification}: 큰 그룹 \(\to\) 세부 분류. 예: 의류 \(\to\) 상의/하의 \(\to\) 셔츠/티셔츠.
\item
  \textbf{희귀 클래스 묶기}: 빈도 1-2회짜리는 “기타”로. long tail 80\%는 어차피 못 배움.
\item
  \textbf{계산 트릭}: hierarchical softmax, sampled softmax (학습 시 일부 클래스만 negative). PyTorch 모델에서 자주, sklearn 기본엔 없음.
\end{enumerate}

K=5 정도는 위 트릭이 필요 없는 작은 규모.

\end{faqBox}

\begin{faqBox}{Q3. (실무) 클래스 불균형이 심한데 \inlinecode{class\_weight}를 어떻게 적용하나요?}

\Needspace{6\baselineskip}
\begin{lstlisting}
model = LogisticRegression(
    class_weight="balanced",
    max_iter=1000,
)
\end{lstlisting}

\inlinecode{balanced}는 빈도 적은 클래스에 더 큰 가중치를 줍니다. 또는 \texttt{class\_weight=\{0:\ 0.5,\ 1:\ 1.0,\ ...\}} 으로 직접 지정 가능. 효과가 미미하면 데이터 레벨 처리(SMOTE 등) 또는 임계값 후처리.

\end{faqBox}

\begin{faqBox}{Q4. (실무) confusion matrix를 시각화하는 추천 방법은?}

\inlinecode{seaborn.heatmap}이 명확합니다.

\Needspace{7\baselineskip}
\begin{lstlisting}
import seaborn as sns
sns.heatmap(cm, annot=True, fmt="d", cmap="Blues",
            xticklabels=[f"{i+1}★" for i in range(5)],
            yticklabels=[f"{i+1}★" for i in range(5)])
plt.xlabel("Predicted"); plt.ylabel("True")
\end{lstlisting}

\textbf{행 정규화}(recall 관점)도 자주 봅니다.

\Needspace{5\baselineskip}
\begin{lstlisting}
cm_norm = cm / cm.sum(axis=1, keepdims=True)
sns.heatmap(cm_norm, annot=True, fmt=".2f", cmap="Blues")
\end{lstlisting}

\end{faqBox}

\begin{faqBox}{Q5. (이론) macro F1과 weighted F1 중 어떤 걸 봐야 하나요?}

\begin{itemize}
\tightlist
\item
  \textbf{macro F1}: 클래스별 F1을 단순 평균. 클래스 불균형이 있으면 소수 클래스 성능을 가리지 않음.
\item
  \textbf{weighted F1}: 클래스 빈도로 가중 평균. 다수 클래스 성능을 더 반영.
\end{itemize}

\textbf{언제 무엇:}

\begin{itemize}
\tightlist
\item
  모든 클래스가 똑같이 중요 \(\to\) macro F1.
\item
  다수 클래스 위주로 평가하고 싶음 \(\to\) weighted F1.
\item
  보고에는 보통 둘 다 함께 적습니다.
\end{itemize}

\end{faqBox}

\begin{faqBox}{Q6. (이론) baseline = \(\log K\)의 의미는 무엇이고 왜 챙겨봐야 하나요?}

CE 식에 균등 분포 \(\hat p_k = 1/K\) 를 대입하면:

\begin{equation}
\label{eq:ch05-02}
L = -\log(1/K) = \log K
\end{equation}

식~\eqref{eq:ch05-02}은 균등 추측 baseline이 \(\log K\)가 되는 이유를 설명합니다.

\textbf{의미}: 모델이 \emph{아무것도 안 학습한 상태} 의 손실. 학습 손실이 baseline보다 크면 “정답이 \emph{아닌} 곳에 자신 있다”는 뜻 --- 모델이 잘못된 방향으로 가고 있습니다.

K=5에서 baseline \(\approx\) 1.609. 학습 도중 loss가 이 값 근처에서 정체되면 “모델이 데이터에서 신호를 못 잡고 있다”는 진단.
\end{faqBox}



\section{오류 실험 (선택)}

같은 5클래스 데이터에서 회귀 모델(\ref{ch:02}장 방식)과 분류 모델(이 장)이 4★를 3★로 예측하는 비율과 4★를 1★로 예측하는 비율의 비를 각각 계산해보세요. 회귀가 더 큰 실수를 더 강하게 회피하는지 확인할 수 있습니다.

\Needspace{11\baselineskip}
\begin{lstlisting}
from sklearn.linear_model import LinearRegression

model_reg = LinearRegression()
model_reg.fit(X_train, y_train.astype(float))
pred_reg = np.clip(np.round(model_reg.predict(X_test)), 0, 4).astype(int)

pred_clf = y_pred

for name, pred in [("회귀 후 round/clip", pred_reg), ("분류", pred_clf)]:
    mask4 = (y_test == 3)
    n_4to3 = ((pred == 2) & mask4).sum()
    n_4to1 = ((pred == 0) & mask4).sum()
    print(
        f"{name}: 4★$\to$3★ {n_4to3}건, 4★$\to$1★ "
        f"{n_4to1}건"
    )
\end{lstlisting}

힌트: 회귀는 큰 실수일수록 손실이 제곱으로 커지므로 4★을 1★로 예측하는 큰 실수가 더 드물어야 합니다. 분류는 둘을 동등하게 처벌하므로 그런 경향이 약할 수 있습니다.



\begin{previewBox}{미리보기: 다음 장}

\textbf{\ref{ch:06}장. 다중 라벨 분류 (Multi-label Classification \& Per-label BCE) --- softmax 합=1 제약을 푼다}

\begin{itemize}
\tightlist
\item
  한 샘플에 \emph{여러} 라벨이 동시에 붙는 multi-label 문제로 확장
\item
  새 데이터: Yelp 리뷰 + \textbf{항목(aspect) 키워드 합성} (food/service/price/ambiance/location 5개)
\item
  softmax 한 개 대신 \textbf{5개 독립 sigmoid} --- 각 라벨이 다른 라벨에 영향받지 않음
\item
  Loss는 CrossEntropyLoss에서 \textbf{per-label \inlinecode{BCEWithLogitsLoss}} 로
\item
  \inlinecode{OneVsRestClassifier(LogisticRegression())} + \inlinecode{MultiLabelBinarizer}로 구현
\end{itemize}
\end{previewBox}